% ============================================================
\section{Cumprimento dos objetivos}\label{sec:objetivos}
% ============================================================

A Tabela~\ref{tab:objetivos} confronta cada objetivo da proposta com o que foi feito e com o
resultado obtido.

% longtable, e não table: a tabela pode começar no fim de uma página e continuar na seguinte.
\begingroup
\small
\renewcommand{\arraystretch}{1.08}
\setlength{\LTpre}{2pt}\setlength{\LTpost}{0pt}
\begin{longtable}{@{}>{\raggedright\arraybackslash}p{3.5cm}
    >{\raggedright\arraybackslash}p{10.7cm} >{\raggedright\arraybackslash}p{1.9cm}@{}}
\caption{Os objetivos da proposta, como foram cumpridos e o que os resultados
mostram.}\label{tab:objetivos}\\
\toprule
\textbf{Objetivo} & \textbf{Como foi cumprido e o que mostrou} & \textbf{Situação} \\
\midrule
\endfirsthead
\multicolumn{3}{@{}l}{\footnotesize\itshape Tabela~\ref{tab:objetivos} (continuação)}\\
\toprule
\textbf{Objetivo} & \textbf{Como foi cumprido e o que mostrou} & \textbf{Situação} \\
\midrule
\endhead
\textbf{Geral}: caracterizar quanto e de que forma a estrutura de similaridade muda entre a
entrada e as camadas inicial, intermediária e final &
Quatro redes k-NN com os mesmos \nVertices{} vértices, comparadas vértice a vértice (troca de
vizinhos, dominância), em escala global (graus, clusterização, distâncias, componentes) e por
comunidades; curvas nos 36 blocos. A mudança é grande, não monótona e depende da categoria do
token (Seção~\ref{sec:discussao}). & cumprido \\
\midrule
\textbf{P1}: quanto e de que forma a organização muda &
Dominância de 0,70 para 0,25 (acaso abaixo de 0,004); Jaccard com a rede lexical de 0,10 no
bloco 36, bem abaixo do piso de ruído; assimetria do grau de entrada de 0 para 4,9; distância
média de 7,6 para 5,3; concordância das comunidades com o token de 0,32 para 0,19
(Seção~\ref{sec:res-p1}). & H1 confirmada \\
\midrule
\textbf{P2}: em que camadas a mudança se concentra &
Transições lexical $\to$ 1 $\to$ 18 $\to$ 36 e curvas nos 36 blocos, por faixa de frequência e
categoria. Continuações e números perdem a identidade até o bloco 4; palavras inteiras, sobretudo
depois do bloco 28; inícios de palavra quase não perdem (Seção~\ref{sec:res-p2}). & H2 confirmada \\
\midrule
\textbf{P3}: comunidades associadas a contextos &
Leiden e NMI $-$ base com oito rótulos, no núcleo. O contexto (artigo, parágrafo, tema) supera
o token nos blocos 12 a 19 e cai depois; o próximo token sobe a partir do bloco 21
(Seção~\ref{sec:res-p3}). & H3 em parte \\
\midrule
\textbf{P4}: separação das ocorrências de um token &
Alvos, controles e multitema com cotas por tema e sentidos anotados. NMI $-$ base comunidade
$\times$ tema de 0,33 nos alvos contra 0,04 nos controles no bloco 36; separação maior pelo
sentido que pelo tema nas 11 palavras (Seção~\ref{sec:res-p4}). & H4 confirmada \\
\bottomrule
\end{longtable}
\endgroup
